% Standalone problem document, generated from the adjacent README.md.
% Compile with XeLaTeX. Mathematical content is maintained in Markdown.
\documentclass[11pt,a4paper]{article}
\usepackage[fontsize=10.5pt]{scrextend}
\usepackage[margin=22mm,headheight=15pt,headsep=9mm,footskip=11mm]{geometry}
\usepackage{fontspec}
\setmainfont{texgyrepagella}[Extension=.otf,UprightFont=*-regular,BoldFont=*-bold,ItalicFont=*-italic,BoldItalicFont=*-bolditalic]
\setsansfont{texgyreheros}[Extension=.otf,UprightFont=*-regular,BoldFont=*-bold,ItalicFont=*-italic,BoldItalicFont=*-bolditalic]
\setmonofont{texgyrecursor}[Extension=.otf,UprightFont=*-regular,BoldFont=*-bold,ItalicFont=*-italic,BoldItalicFont=*-bolditalic,Scale=MatchLowercase]
\usepackage{amsmath,amssymb}
\usepackage{unicode-math}
\setmathfont{texgyrepagella-math.otf}
\usepackage{xcolor,microtype,enumitem,needspace,fancyhdr}
\usepackage{bookmark}
\definecolor{ink}{HTML}{17344B}
\definecolor{accent}{HTML}{197D80}
\definecolor{muted}{HTML}{596773}
\hypersetup{colorlinks=true,linkcolor=accent,urlcolor=accent,pdftitle={NM-04 - The
Rowland--Wu polynomial identity for Sinkhorn
limits},pdfauthor={Open Problems in Numerical Linear Algebra}}
\urlstyle{same}
\setlength{\parindent}{0pt}
\setlength{\parskip}{5pt plus 1pt minus 1pt}
\linespread{1.04}
\setlength{\emergencystretch}{3em}
\widowpenalty=10000
\clubpenalty=10000
\displaywidowpenalty=10000
\allowdisplaybreaks[1]
\setlist{nosep,leftmargin=1.5em}
\providecommand{\tightlist}{\setlength{\itemsep}{0pt}\setlength{\parskip}{0pt}}
\providecommand{\pandocbounded}[1]{#1}
\pagestyle{fancy}
\fancyhf{}
\fancyhead[L]{\sffamily\footnotesize\color{muted}OPEN PROBLEMS IN NUMERICAL LINEAR ALGEBRA}
\fancyhead[R]{\sffamily\bfseries\footnotesize\color{accent}NM-04}
\fancyfoot[L]{\parbox[b]{0.9\textwidth}{\sffamily\fontsize{8}{10}\selectfont\color{muted}Editorial ratings; a bounded search cannot certify that a problem remains unsolved.\\Literature check: 2026-09-10}}
\fancyfoot[R]{\sffamily\footnotesize\color{muted}\thepage}
\renewcommand{\headrulewidth}{0.3pt}
\renewcommand{\headrule}{\hbox to\headwidth{\color{accent}\leaders\hrule height \headrulewidth\hfill}}
\makeatletter
\renewcommand{\section}{\Needspace{5\baselineskip}\@startsection{section}{1}{0pt}{12pt plus 2pt minus 2pt}{4pt}{\sffamily\bfseries\large\color{ink}}}
\renewcommand{\subsection}{\Needspace{4\baselineskip}\@startsection{subsection}{2}{0pt}{9pt plus 1pt minus 1pt}{3pt}{\sffamily\bfseries\normalsize\color{ink}}}
\makeatother
\setcounter{secnumdepth}{0}
\begin{document}
{\sffamily\small\color{muted}Nonnegative and positive
factorizations\par}
\vspace{3pt}
{\raggedright\hyphenpenalty=10000\exhyphenpenalty=10000\sffamily\bfseries\fontsize{21}{24}\selectfont\color{ink}The
Rowland--Wu polynomial identity for Sinkhorn limits\par}
\vspace{7pt}
{\sffamily\small\textbf{Difficulty:} challenging\quad\textbf{Importance:} interesting
to the community\par}
{\sffamily\small\bfseries\color{accent}Status: Partially resolved\par}
\vspace{4pt}
{\color{accent}\hrule height 0.8pt}
\vspace{6pt}
\textbf{Topic:} Exact matrix scaling and balancing.

\textbf{Rating rationale:} Challenging because the full coefficient
formula requires a general algebraic-combinatorial identity; community
importance concerns exact formulas for a widely used matrix-scaling
limit.

\subsection{Problem statement}\label{problem-statement}

Let \(m,n\ge1\) and \(A\in\mathbb R_{>0}^{m\times n}\). Let
\(\operatorname{Sink}(A)\) be its Sinkhorn limit with row sums \(1\) and
column sums \(m/n\): equivalently, the unique matrix \(D_1AD_2\) with
these sums, where \(D_1,D_2\) are positive diagonal matrices. Set
\(x=\operatorname{Sink}(A)_{11}\).

Define

\[
\mathcal D=\{(R,C):R\subseteq\{2,\ldots,m\},\ C\subseteq\{2,\ldots,n\},\ |R|=|C|\}.
\]

Submatrix indices are in increasing order; empty determinants and
products equal \(1\). Put

\[
\Delta(R,C)=\det A_{\{1\}\cup R,\{1\}\cup C},\qquad
\Gamma(R,C)=a_{11}\det A_{R,C},
\]

\[
M(\mathcal S)=\prod_{(R,C)\in\mathcal S}\Delta(R,C)
\prod_{(R,C)\in\mathcal D\setminus\mathcal S}\Gamma(R,C).
\]

For a finite set \(U\) of integers and \(s\in U\), let \(p_s(U)\) be the
position of \(s\) in the increasing list of \(U\), starting at \(1\).
For each \(\mathcal S=\{(R_i,C_i):1\le i\le k\}\subseteq\mathcal D\),
define \(H_{\mathcal S}\in\mathbb Z^{k\times k}\) by

\[
(H_{\mathcal S})_{ii}=|R_i|(m+n)-mn,
\qquad
\tau_{ij}=(R_i\setminus R_j,R_j\setminus R_i,C_i\setminus C_j,C_j\setminus C_i),
\]

and, for \(i\ne j\),

\[
(H_{\mathcal S})_{ij}=\begin{cases}
(-1)^{p_s(R_j)+p_t(C_j)}m,&\tau_{ij}=(\varnothing,\{s\},\varnothing,\{t\}),\\
(-1)^{p_s(R_i)+p_t(C_i)+1}n,&\tau_{ij}=(\{s\},\varnothing,\{t\},\varnothing),\\
(-1)^{p_s(C_i)+p_t(C_j)}m,&\tau_{ij}=(\varnothing,\varnothing,\{s\},\{t\}),\\
(-1)^{p_s(R_i)+p_t(R_j)}n,&\tau_{ij}=(\{s\},\{t\},\varnothing,\varnothing),\\
0,&\text{otherwise}.
\end{cases}
\]

The determinant is independent of the ordering chosen for
\(\mathcal S\). Is the following identity valid for every such
\(A,m,n\)?

\[
\sum_{\mathcal S\subseteq\mathcal D}
\det(m^{-1}H_{\mathcal S})M(\mathcal S)x^{|\mathcal S|}=0.
\]

\subsection{Why it matters}\label{why-it-matters}

This would give an explicit algebraic relation for entries of a
ubiquitous iterative matrix scaling limit, including a structured
formula for every coefficient.

\subsection{References}\label{references}

\begin{itemize}
\tightlist
\item
  E. Rowland and J. Wu, \href{https://arxiv.org/abs/2409.02789}{The
  entries of the Sinkhorn limit of an \(m\times n\) matrix}, arXiv v2
  (2025-05-25), notation on pp.~3--4 and Conjecture 2. The matrix
  denoted \(H_{\mathcal S}\) here is their
  \(\operatorname{adj}_{\mathcal S}(m,n)\).
\item
  E. Rowland,
  \href{https://ericrowland.github.io/talks/Combinatorial_structure_behind_Sinkhorn_limits_SIAM.pdf}{Combinatorial
  structure behind Sinkhorn limits}, SIAM talk, 2025-07-07, slide 13:
  the coefficient formula remains conjectural after the degree bound was
  proved.
\item
  M. C. Fang, \href{https://arxiv.org/abs/2506.06338}{Closed Form of a
  Generalized Sinkhorn Limit}, 2025, abstract and the general algebraic
  degree bound. This settles the degree-bound consequence, rather than
  the displayed coefficient formula.
\end{itemize}

\subsection{Status check --- 2026-09-10}\label{status-check-2026-09-10}

Rechecked \href{https://arxiv.org/pdf/2409.02789}{Rowland--Wu v2, §2 and
Conjecture 2}, including the rectangular scaling and four coefficient
signs, and \href{https://arxiv.org/abs/2506.06338}{Fang's degree-bound
result}. The 3×3 coefficient formula is proved, while the determinant
formula in arbitrary dimensions remains conjectural. Searches for
subsequent Rowland--Wu identity proofs found no full resolution. Fang's
general degree bound alone does not establish the displayed
coefficients.
\end{document}
